\documentclass[11pt]{article}
\usepackage[a4paper,margin=28mm]{geometry}
\usepackage{amsmath,amssymb,amsthm,hyperref}
\hypersetup{colorlinks=true,urlcolor=blue,citecolor=blue,linkcolor=blue}
\newtheorem{theorem}{Theorem}
\newtheorem{lemma}[theorem]{Lemma}
\newtheorem{corollary}[theorem]{Corollary}
\theoremstyle{remark}\newtheorem{remark}[theorem]{Remark}
\newcommand{\Q}{\mathbb Q}\newcommand{\Z}{\mathbb Z}
\newcommand{\F}{\mathbb F}\newcommand{\T}{\mathbb T}
\newcommand{\ee}{\mathbf e}
\DeclareMathOperator{\Cot}{Cot}\DeclareMathOperator{\Lie}{Lie}
\DeclareMathOperator{\Ann}{Ann}\DeclareMathOperator{\Hom}{Hom}
\DeclareMathOperator{\Tr}{Tr}
\title{A dyadic bound for torsion primes}
\author{}
\date{27 October 2026}
\begin{document}
\maketitle
\begin{abstract}
We prove the bound
$p\leq\max\{(2^{d/2}+1)^2,65(4d)^6\}$ for prime-order torsion on
elliptic curves over number fields of degree $d\geq3$.
In particular, $p\leq(2^{d/2}+1)^2$ for every $d\geq53$.
The argument combines Parent's modular-symbol estimate with Raynaud's
bound on the failure of exactness of cotangent lattices at $2$.
Doubling a trace divisor moves it into a residue neighbourhood modulo $4$,
where the resulting map has a unique zero.
\end{abstract}

Let $S(d)$ be the set of primes occurring as orders of rational points
on elliptic curves over number fields of degree $d$.
Oesterl\'e's bound is
$p\leq(3^{d/2}+1)^2$ for $p\in S(d)$; see \cite[Section~6]{DKSS}.
We prove the following statement.

\begin{theorem}\label{thm:main}
Let $d\geq3$, let $K/\Q$ have degree $d$, and let $E/K$ be an elliptic
curve with a $K$-rational point of prime order $p$. Then
\[
 p\leq\max\bigl\{(2^{d/2}+1)^2,\;65(4d)^6\bigr\}.
\]
Consequently, for every $d\geq53$, every prime divisor of
$\#E(K)_{\rm tors}$ satisfies
\[
 p\leq(2^{d/2}+1)^2.
\]
\end{theorem}

We first record the local observations needed in the proof.
All cotangent spaces below are relative to the indicated base ring.

\begin{lemma}\label{lem:local}
Let $f$ be a morphism between smooth $\Z_2$-schemes, let $x_0$ and $y_0$
be sections with $f(x_0)=y_0$, and choose coordinates
$x_1,\ldots,x_m$ for the formal neighbourhood of $x_0$.
Suppose that the cotangent image of $f$ at this section contains
$2\Z_2^m$. Then the only point with all $x_i\in4\Z_2$ mapping to
$y_0$ is $x_0$.
\end{lemma}
\begin{proof}
Choose integral formal functions on the target, vanishing at $y_0$,
whose pullbacks are
\[
 G_i(x)=2x_i+Q_i(x),\qquad 1\leq i\leq m,
\]
where each $Q_i$ belongs to $(x_1,\ldots,x_m)^2\Z_2[[x_1,\ldots,x_m]]$.
Suppose that $x\ne0$ and put $k=\min_i v_2(x_i)\geq2$.
For an index attaining this minimum, the linear term of $G_i(x)$ has
valuation $k+1$, whereas $Q_i(x)$ has valuation at least $2k>k+1$.
Hence $G_i(x)\ne0$.
\end{proof}

\begin{lemma}\label{lem:double}
Let $X/\Z_2$ be a smooth curve with a section $c$, and let
$D\in X^{(d)}(\Z_2)$ reduce to $dc$ modulo $2$.
Then $2D\in X^{(2d)}(\Z_2)$ agrees with $2dc$ modulo $4$ in the
formal neighbourhood of $2dc$.
\end{lemma}
\begin{proof}
Choose a parameter $q$ along $c$. At $dc$, the elementary symmetric
functions of the local parameters give coordinates on $X^{(d)}$.
Encode their values by
\[
 E_D(t)=\prod_{j=1}^d(1-q_jt)=1+2B(t),\qquad B(t)\in\Z_2[t].
\]
This expression is defined over $\Z_2$ even when the individual $q_j$
are defined over ramified extensions or occur with repetitions.
The corresponding polynomial for the doubled divisor is
\[
 E_{2D}(t)=E_D(t)^2=1+4\bigl(B(t)+B(t)^2\bigr).
\]
Thus every nonconstant elementary symmetric coordinate of $2D$ is
divisible by $4$.
\end{proof}

For an abelian scheme $A/\Z_2$, the torsion in the kernel of
$A(\Q_2)\to A(\F_2)$ is killed by $2$.
Indeed, multiplication by $2$ sends the formal neighbourhood
with coordinates in $2\Z_2$ into the neighbourhood with coordinates
in $4\Z_2$. The latter is torsion free, either by the formal logarithm
or the usual valuation argument for the multiplication series.
We use this fact for the winding quotient.

Let $p$ be an odd prime, put $J=J_0(p)$, and let $\T$ be the integral
Hecke algebra acting on $J$. Let $\ee$ be the winding element and put
\[
 I=\Ann_{\T}(\ee),\qquad A=J/IJ,\qquad \pi:J\longrightarrow A.
\]
The ideal $I$ is saturated and $\pi$ has connected kernel.
The winding quotient $A(\Q)$ is finite; see
\cite[Section~3.8]{Parent}.
Since $p$ is odd, all these abelian varieties have good reduction at $2$.
Write
\[
 S=H^0(X_0(p)_{\Z_2},\Omega^1),\qquad
 M=S[I],\qquad L=\pi^*\Cot_0(A_{\Z_2})\subseteq S.
\]

\begin{lemma}\label{lem:lattice}
There are inclusions
\[
 2M\subseteq L\subseteq M.
\]
Moreover, the coefficient pairing gives an identification
\[
 M=\Hom_{\Z_2}\bigl((\T/I)\otimes\Z_2,\Z_2\bigr),
 \qquad f\longmapsto\bigl(T_n\longmapsto a_n(f)\bigr).
\]
\end{lemma}
\begin{proof}
Put $B=\ker\pi$. Raynaud's theorem
\cite[Appendix~A, Theorem~A.1]{AU}, applied to the exact sequence
$0\to B\to J\to A\to0$ over $\Q_2$, gives
\[
 0\longrightarrow\Lie B_{\Z_2}\longrightarrow\Lie J_{\Z_2}
 \longrightarrow\Lie A_{\Z_2}\longrightarrow(\Z/2\Z)^r
 \longrightarrow0
\]
for some $r$. Here the ramification index is $1=2-1$, and the middle
term has good reduction. Dualizing shows that the saturation of $L$
in $S$ contains $L$ with quotient killed by $2$; compare
\cite[Section~4.4, proof of Theorem~3.10]{ARS}.
The $\Q_2$-span of this saturation is the subspace annihilated
by $I$. Hence its integral lattice is $M$, proving the inclusions.

The integral coefficient pairing between $\T\otimes\Z_2$ and $S$
is perfect by \cite[Lemma~2.1]{AU} and the integral
$q$-expansion identification at primes not dividing the level.
The submodule annihilated by $I$ is consequently the dual of
$(\T/I)\otimes\Z_2$. No Gorenstein hypothesis is needed.
\end{proof}

We next use the part of Parent's argument which concerns modular symbols.
For $m\geq3$, his combinatorial estimate implies that
\[
 T_i\{0,\infty\},\qquad 1\leq i\leq2m,
\]
are linearly independent in relative homology with coefficients in any
field whenever $p>65(2m)^6$; see \cite[Proposition~1.9]{Parent}.
The corresponding result is Proposition~1.7 in the 1996 preprint.
The following passage to winding symbols also works in characteristic $2$.

\begin{lemma}\label{lem:jets}
If $m\geq3$ and $p>65(2m)^6$, then the map
\[
 M\longrightarrow\Z_2^m,\qquad
 f\longmapsto(a_1(f),\ldots,a_m(f))
\]
is surjective.
Consequently, the cotangent image at $m\infty$ of
\[
 f_m:X_0(p)^{(m)}\longrightarrow A,\qquad
 D\longmapsto\pi[D-m\infty],
\]
contains $2\Z_2^m$.
\end{lemma}
\begin{proof}
We show that $T_1\ee,\ldots,T_m\ee$ are independent modulo $2\T\ee$.
Put $\gamma=\{0,\infty\}$ and $h=T_2-3$.
Both cusps are multiplied by $3$ under $T_2$, so $h\gamma$ lies in
integral absolute homology. The period definition of the winding element
gives $h\gamma=h\ee$.
In particular, $h\T\ee$ lies in integral absolute homology.

Suppose that
\[
 \sum_{i=1}^m\epsilon_iT_i\ee\in2\T\ee,
 \qquad \epsilon_i\in\{0,1\},
\]
and let $k$ be the largest index with $\epsilon_k=1$.
Multiplication by $h$ gives a dependence modulo $2$ between
$T_i\gamma$ for $1\leq i\leq2k$, whose coefficient at $T_{2k}\gamma$
is $1$. Indeed,
\[
 (T_2-3)T_i=T_{2i}+2T_{i/2}-3T_i,
\]
where $T_{i/2}=0$ for odd $i$.
This contradicts Parent's estimate. This argument takes place in integral
homology and uses no division by $2$.

Since $\T\ee\simeq\T/I$, independence modulo $2$ makes the submodule
generated by the first $m$ operators primitive over $\Z_2$.
Dualizing and applying Lemma~\ref{lem:lattice} proves the coefficient
surjectivity for $M$.
The cotangent map of $f_m$ in elementary symmetric coordinates is
\[
 \omega\longmapsto
 \bigl(a_1(\pi^*\omega),-a_2(\pi^*\omega),\ldots,
 (-1)^{m-1}a_m(\pi^*\omega)\bigr);
\]
see \cite[Section~4.12]{Parent}.
The inclusion $2M\subseteq L$ proves the last assertion.
\end{proof}

\begin{proof}[Proof of Theorem~\ref{thm:main}]
Suppose that $p$ exceeds both terms in the maximum, and let
$P\in E(K)$ have order $p$.
At a place above $2$ with residue degree $f\leq d$, good reduction
would imply
\[
 p\leq\#\widetilde E(\F_{2^f})\leq(2^{f/2}+1)^2
 \leq(2^{d/2}+1)^2.
\]
Additive reduction is impossible since its odd torsion has order at most $3$.
Nonsplit multiplicative reduction is also impossible: its odd torsion
injects into a torus of order $2^f+1$.
Thus $E$ has split multiplicative reduction at every place above $2$.
The reduction of $P$ has order $p$ in the component group, since
$p>2^f-1$.

Let $x\in X_0(p)(K)$ be the image of $(E,\langle P\rangle)$ followed
by the Atkin--Lehner involution. The component-group assertion says
that $x$ specializes to $\infty$ at every place above $2$.
Hence
\[
 D=\Tr_{K/\Q}(x)\in X_0(p)^{(d)}(\Q)
\]
reduces to $d\infty$ modulo $2$. This statement allows repeated
points when $x$ is defined over a proper subfield of $K$.

The point $f_d(D)\in A(\Q)$ is torsion and reduces to zero at $2$.
The reduction-kernel observation therefore gives
\[
 2f_d(D)=0,\qquad f_{2d}(2D)=0.
\]
By Lemma~\ref{lem:double}, $2D$ lies in the residue neighbourhood
of $2d\infty$ modulo $4$.
Put $m=2d$. Our assumption $p>65(4d)^6=65(2m)^6$ and
Lemma~\ref{lem:jets} allow us to apply Lemma~\ref{lem:local} to $f_m$.
It follows that $2D=2d\infty$, contradicting the noncuspidal support of $D$.

Finally,
\[
 65(4\cdot53)^6=5901039506984960<9007199254740992=2^{53}.
\]
The ratio $2^d/d^6$ is increasing for integers $d\geq53$, since
$2(53/54)^6>1$ and $d/(d+1)$ is increasing.
Thus $65(4d)^6<2^d<(2^{d/2}+1)^2$ for every $d\geq53$.
\end{proof}

\begin{remark}
The theorem bounds torsion primes. It does not assert the same bound
for the cardinality of the full torsion group or for arbitrary prime powers.
The proof uses the exponent-two reduction-kernel statement and Raynaud's
exponent-two cotangent defect. It does not require a uniform formal
immersion theorem modulo $2$, nor a uniform rank assertion for $J_0(p)$.
\end{remark}

\begin{thebibliography}{9}
\bibitem{AU}
A. Abbes and E. Ullmo,
\emph{\`A propos de la conjecture de Manin pour les courbes elliptiques modulaires},
Compositio Math. \textbf{103} (1996), 269--286.
\url{https://www.numdam.org/item/CM_1996__103_3_269_0/}.

\bibitem{ARS}
A. Agashe, K. Ribet and W. A. Stein,
\emph{The Manin constant},
Pure Appl. Math. Q. \textbf{2} (2006), 617--636.
\url{https://wstein.org/papers/ars-manin/html/node12.html}.

\bibitem{DKSS}
M. Derickx, S. Kamienny, W. Stein and M. Stoll,
\emph{Torsion points on elliptic curves over number fields of small degree},
Algebra Number Theory \textbf{17} (2023), 267--308.
\url{https://arxiv.org/abs/1707.00364}.

\bibitem{Parent}
P. Parent,
\emph{Bornes effectives pour la torsion des courbes elliptiques sur les corps de nombres},
J. Reine Angew. Math. \textbf{506} (1999), 85--116.
Preprint: \url{https://arxiv.org/abs/alg-geom/9611022}.
\end{thebibliography}
\end{document}
